\input{configTD}

% ── Poly à trous : commandes spécifiques ──────────────────────────

% Bloc à compléter en cours.
% Mode corrigé : contenu sur fond coloré (mdframed).
% Mode énoncé  : cadre gris vide de hauteur #1 (défaut 3 cm).
\newcommand{\atrou}[2][3cm]{%
  \ifthenelse{\boolean{showSolutions}}{%
    \par\vspace{0.5em}%
    \begin{mdframed}[backgroundcolor=ThemeLight, linewidth=0pt,
      skipabove=0pt, skipbelow=0pt,
      innertopmargin=8pt, innerbottommargin=8pt]%
    #2%
    \end{mdframed}%
    \par\vspace{0.5em}%
  }{%
    \par\vspace{0.5em}%
    \noindent\fbox{\begin{minipage}[c][#1]%
      {\dimexpr\textwidth-2\fboxsep-2\fboxrule\relax}%
      \color{white}\rule{0pt}{0pt}%
    \end{minipage}}%
    \par\vspace{0.5em}%
  }%
}

% Valeur dans un tableau : visible en corrigé, fantôme en énoncé.
\newcommand{\val}[1]{%
  \ifthenelse{\boolean{showSolutions}}{#1}{\phantom{#1}}%
}

% Espace pour un croquis.
\newcommand{\esquisse}[1][5.5cm]{%
  \begin{center}%
  \ifthenelse{\boolean{showSolutions}}{%
    \fbox{\begin{minipage}[c][#1]{0.85\textwidth}%
      \centering\vfill
      \textit{\small Croquis réalisé en cours :
      courbe exacte, tangente au point $(t_n,T_n)$,
      pas d'Euler, écart avec la solution.}%
      \vfill
    \end{minipage}}%
  }{%
    \fbox{\begin{minipage}[c][#1]{0.85\textwidth}%
      \hfill\vfill\hfill
    \end{minipage}}%
  }%
  \end{center}%
}

\newcommand{\degC}{\text{\textdegree C}}

\title{\sffamily\bfseries Analyse Numérique --- CM1\\[0.3em]
  \Large De l'erreur numérique à la méthode d'Euler}
\author{}
\date{17 mars 2026}

\begin{document}
\sffamily
\maketitle
\thispagestyle{empty}

% ====================================================================
\section*{1. Le problème}
% ====================================================================

Partons d'une situation très simple,
qui ne demande a priori ni formule ni ordinateur :

\begin{center}
\begin{mdframed}[backgroundcolor=ThemeLight, linecolor=Theme,
  linewidth=1.5pt, innertopmargin=10pt, innerbottommargin=10pt]
\centering\large\itshape
\og Voici un café trop chaud.
Combien de temps faut-il attendre pour pouvoir le boire ?\fg{}
\end{mdframed}
\end{center}

Pour y répondre, il faut d'abord se poser les bonnes questions.

\medskip

\textbf{Questions préliminaires.}
\begin{itemize}[leftmargin=*]
  \item Quelle grandeur varie au cours du temps ?
  \item Comment modéliser le système ?
\end{itemize}

\atrou[3.6cm]{%
\begin{itemize}[leftmargin=*]
  \item La température $T(t)$ varie au cours du temps.
  \item On peut modéliser le système par une équation différentielle : Plus la température du café est élevée, plus il refroidit vite.
  L'équation serait de type $T'(t) = -kT(t)$.

  A l'équilibre, la température du café est égale à la température de la pièce:
  L'équation devient $T'(t) = -k(T(t)-T_{\mathrm{amb}})$.
\end{itemize}
}

% ====================================================================
\section*{2. Modélisation --- Bilan thermique}
% ====================================================================

Pour répondre à la question, il faut un modèle :
une équation qui décrit comment la température évolue.

\medskip

On considère un café de température $T(t)$ dans une pièce à température
ambiante constante $T_{\mathrm{amb}}$.
En supposant la température uniforme dans la tasse
(modèle à paramètres localisés), le bilan d'énergie s'écrit :
\[
  m\,c_p\,\frac{dT}{dt}
  = -h_{\mathrm{eff}}\,A\,\bigl(T(t) - T_{\mathrm{amb}}\bigr).
\]

\begin{center}
\renewcommand{\arraystretch}{1.2}
\begin{tabular}{llcl}
  \hline
  \textbf{Grandeur} & \textbf{Symbole}
    & \textbf{Valeur} & \textbf{Unité} \\
  \hline
  Masse de café
    & $m$ & $0{,}25$ & kg \\
  Capacité thermique massique
    & $c_p$ & $4180$ & $\mathrm{J\,kg^{-1}\,K^{-1}}$ \\
  Coefficient d'échange global
    & $h_{\mathrm{eff}}$ & $\approx 50$
    & $\mathrm{W\,m^{-2}\,K^{-1}}$ \\
  Surface d'échange de la tasse
    & $A$ & $\approx 0{,}035$ & $\mathrm{m}^{2}$ \\
  \hline
\end{tabular}
\end{center}

\smallskip

\textit{Le coefficient $h_{\mathrm{eff}}$ regroupe convection naturelle,
rayonnement et évaporation (dominante pour un café chaud non couvert).}

\medskip

\textbf{Forme simplifiée.} En divisant par $m\,c_p$, on peut simplifier l'équation en :

\[
  \frac{dT}{dt} = -k\bigl(T(t) - T_{\mathrm{amb}}\bigr),
  \qquad k = \frac{h_{\mathrm{eff}}\,A}{m\,c_p}\,.
\]

\textbf{Application numérique :}
\[
  k = \frac{50 \times 0{,}035}{0{,}25 \times 4180}
    = \frac{1{,}75}{1045}
    \approx 0{,}0017\;\mathrm{s^{-1}}
    \approx 0{,}10\;\mathrm{min^{-1}}.
\]

On vient d'obtenir une EDO à coefficients constants, on peut la résoudre de manière exacte.



% ====================================================================
\section*{3. A quoi sert l'analyse numérique ?}
% ====================================================================

But du module : faire résoudre par un ordinateur des problèmes issues de différentes modélisations. 

Dans l'exemple précédent, on peut tout résoudre à la main, mais ce modèle est volontairement simplifié :

\medskip

\textbf{Complexification du café.}

Si on veut pousser un peu notre modélisation, $h_{\mathrm{eff}}$ n'est pas constant :
il dépend de la température $T$ du café.
Il regroupe trois contributions, chacune non linéaire :

\begin{itemize}[leftmargin=*]
  \item \textbf{Convection naturelle} :
    $h_{\mathrm{conv}}$ est proportionnel à $(T - T_{\mathrm{amb}})^{1/4}$
    \qquad (corrélation de type Churchill-Chu).
  \item \textbf{Rayonnement} :
    $h_{\mathrm{rad}}
    = \varepsilon\,\sigma\,(T^2+T_{\mathrm{amb}}^2)\,(T+T_{\mathrm{amb}})$
    \qquad ($T$ en Kelvin).
  \item \textbf{Évaporation} (terme dominant) :
    le flux évaporatif dépend de la pression
    de vapeur saturante,
    $$p_{\mathrm{sat}}(T)
    \approx 611 \times
    \exp\!\Bigl(\dfrac{17{,}27\,T_C}{T_C+237}\Bigr)
    \;[Pa],$$
    qui croît quasi exponentiellement avec $T$.
\end{itemize}

\smallskip

À $90\;\degC$, l'évaporation représente
${\sim}\,80\,\%$ du transfert total
et $h_{\mathrm{eff}} \approx 100\;\mathrm{W\,m^{-2}\,K^{-1}}$.

À $40\;\degC$, les trois contributions sont comparables
et $h_{\mathrm{eff}} \approx 25\;\mathrm{W\,m^{-2}\,K^{-1}}$.

Le coefficient varie donc d'un facteur $\sim\!4$
entre un café brûlant et un café tiède.

\smallskip

L'équation devient :
\[
  m\,c_p\,\frac{dT}{dt}
  = -\underbrace{\bigl[h_{\mathrm{conv}}(T)
    + h_{\mathrm{rad}}(T)
    + h_{\mathrm{evap}}(T)\bigr]}_{h_{\mathrm{eff}}(T)}
  \,A\,\bigl(T(t) - T_{\mathrm{amb}}\bigr).
\]
Puissance $1/4$, produit de carrés, exponentielle :
chaque terme est non linéaire en $T$.
On ne sait plus résoudre cette équation par une formule.

\medskip

Ce phénomène n'est pas propre au café.
Voici deux autres situations d'ingénierie
qui posent le même problème.

\medskip

\textbf{Vidange d'un réservoir} (loi de Torricelli).
La hauteur d'eau $h(t)$ vérifie :
\[
  \frac{dh}{dt} = -\alpha\,\sqrt{h(t)},
  \qquad h(t) \geq 0.
\]
La non-linéarité en $\sqrt{h}$ change la dynamique :
le réservoir se vide de plus en plus lentement.
Si en plus la section du réservoir varie avec $h$,
aucune formule simple n'est disponible.

\medskip

\textbf{Pendule simple} (grandes oscillations).
L'angle $\theta(t)$ vérifie :
\[
  \theta''(t)
  = -\frac{g}{L}\sin\,\theta(t).
\]
Pour de petites oscillations,
$\sin(\theta)\approx\theta$ et on retrouve
un oscillateur harmonique que l'on sait résoudre.
Mais pour de grandes amplitudes,
le $\sin$ empêche toute solution en termes
de fonctions élémentaires.

\medskip

\textbf{Qu'ont en commun ces trois modèles ?}

\atrou[3.6cm]{%
Dans chaque cas, on connaît une \textbf{loi d'évolution locale}
(une dérivée, un taux de variation),
mais la solution globale n'est pas accessible par une formule.

C'est exactement là que l'analyse numérique intervient :
elle fournit une méthode de calcul approché,
applicable à tous ces problèmes.
}

% ====================================================================
\section*{4. Implémentation dans un ordinateur}
% ====================================================================

Les ordinateurs savent seulement réaliser des sommes et des produits. Pas de calcul d'intégrales, pas de dérivée, pas de fonctions trigonométriques ou de racine carrée. 

Comment les ordinateurs/calculatrices modernes calculent des racines carrées ou des cosinus ? 
\atrou[3.6cm]{%
  La réponse est presque la même pour les deux types de calculs, on stocke des valeurs qu'on a calculées à l'avance avec des méthodes numériques (méthodes de Newton, séries de Taylor, etc.). Si on a besoin d'une valeur intermédiaire, on réalise une interpolation (voir la suite du cours) pour obtenir une valeur approchée.

}


Que nous dit exactement l'équation différentielle ? 

\medskip

Revenons au modèle simple du café :
$\displaystyle\frac{dT}{dt}
= -k\bigl(T(t) - T_{\mathrm{amb}}\bigr).$

\begin{itemize}[leftmargin=*]
  \item Si le café est beaucoup plus chaud que la pièce,
        la température baisse-t-elle vite ou lentement ?
  \item Quand le café approche de $T_{\mathrm{amb}}$,
        que devient la vitesse de refroidissement ?
  \item Si à un instant $t$ on connaît $T(t)$,
        l'équation donne-t-elle directement $T(t+10)$,
        ou seulement la manière dont $T$ évolue à cet instant ?
\end{itemize}

\atrou[4cm]{%
\textbf{Conclusion :}
l'EDO donne une \textbf{information locale} ---
la pente (taux de variation) à l'instant courant.
Elle ne donne pas directement la valeur future.
}

La question qui se pose alors naturellement est :

\begin{center}
\begin{mdframed}[backgroundcolor=ThemeLight, linecolor=Theme,
  linewidth=1pt, innertopmargin=8pt, innerbottommargin=8pt]
\centering\itshape
\og Si je connais la valeur actuelle et la pente actuelle,
comment approcher la valeur un peu plus tard\,?\fg{}
\end{mdframed}
\end{center}

% ====================================================================
\section*{5. Méthode d'Euler explicite}
% ====================================================================

L'idée est simple : si l'on connaît la valeur actuelle
et la pente actuelle, on peut estimer la valeur
un petit moment plus tard en suivant cette pente.
C'est le principe de la méthode d'Euler.

\medskip

On note $T_n \approx T(t_n)$ la valeur approchée au temps
$t_n = t_0 + n\,\Delta t$.

\medskip

\textbf{Construction sur le café.}

\atrou[6.4cm]{%
La pente à l'instant $t_n$ vaut
$f(t_n,T_n) = -k\,(T_n - T_{\mathrm{amb}})$.

En supposant cette pente presque constante sur
$[t_n,\,t_n+\Delta t]$ :
\[
  \boxed{T_{n+1} = T_n
    + \Delta t\left[-k\bigl(T_n - T_{\mathrm{amb}}\bigr)\right].}
\]
}

\medskip

\textbf{Formule générale.}
Pour $y'=f(t,y)$ avec $y(t_0)=y_0$ :

\atrou[4cm]{%
\[
  \boxed{y_{n+1} = y_n + \Delta t\,f(t_n,\,y_n).}
\]

\begin{itemize}
  \item $y_n$ : valeur approchée au temps $t_n$ ;
  \item $\Delta t$ : pas de temps ;
  \item $f(t_n,y_n)$ : pente fournie par le modèle.
\end{itemize}
}

\medskip

\textbf{Interprétation géométrique}
--- tracer la courbe exacte, la tangente au point
$(t_n,T_n)$ et le pas d'Euler :

\esquisse[10cm]

% ====================================================================
\section*{6. Activité --- Euler à la main}
% ====================================================================

On va maintenant appliquer le schéma d'Euler sur le café
avec des valeurs numériques concrètes,
pour voir la méthode en action et comparer
deux pas de temps différents.

\medskip

\textbf{Paramètres :}\;
$T_0 = 90\;\degC$,\;
$T_{\mathrm{amb}} = 20\;\degC$,\;
$k = 0{,}1\;\mathrm{min^{-1}}$.
\[
  T_{n+1} = T_n - 0{,}1\;\Delta t\;(T_n - 20).
\]

\medskip

\textbf{Cas $\boldsymbol{\Delta t = 2}$~min}
\quad(\,calculer $T_1$ à $T_5$,
soit $10$ min de simulation\,) :

\begin{center}
\renewcommand{\arraystretch}{1.6}
\begin{tabular}{|c|c|c|c|c|}
  \hline
  $n$ & $t_n\;\text{(min)}$ & $T_n\;(\degC)$
      & Pente $-0{,}1(T_n-20)$ & $T_{n+1}\;(\degC)$ \\
  \hline
  0 & 0 & 90
    & \val{$-7{,}0$} & \val{$76$} \\
  \hline
  1 & 2 & \val{$76$}
    & \val{$-5{,}6$} & \val{$64{,}8$} \\
  \hline
  2 & 4 & \val{$64{,}8$}
    & \val{$-4{,}48$} & \val{$55{,}84$} \\
  \hline
  3 & 6 & \val{$55{,}84$}
    & \val{$-3{,}584$} & \val{$48{,}67$} \\
  \hline
  4 & 8 & \val{$48{,}67$}
    & \val{$-2{,}867$} & \val{$42{,}94$} \\
  \hline
\end{tabular}
\end{center}

\textbf{Questions.}
\begin{itemize}[leftmargin=*]
  \item La température calculée baisse-t-elle bien
        à chaque étape ?
  \item La correction ajoutée diminue-t-elle d'un pas
        à l'autre ? Pourquoi est-ce cohérent avec le modèle ?
  \item L'évolution semble-t-elle physiquement crédible ?
\end{itemize}

\medskip

\textbf{Cas $\boldsymbol{\Delta t = 5}$~min}
\quad(\,calculer $T_1$ et $T_2$\,) :

\begin{center}
\renewcommand{\arraystretch}{1.6}
\begin{tabular}{|c|c|c|c|c|}
  \hline
  $n$ & $t_n\;\text{(min)}$ & $T_n\;(\degC)$
      & Pente $-0{,}1(T_n-20)$ & $T_{n+1}\;(\degC)$ \\
  \hline
  0 & 0 & 90
    & \val{$-7{,}0$} & \val{$55$} \\
  \hline
  1 & 5 & \val{$55$}
    & \val{$-3{,}5$} & \val{$37{,}5$} \\
  \hline
\end{tabular}
\end{center}

\medskip

\textbf{Solution exacte} à $t=10$ :
$T(10) = 20 + 70\,e^{-1} \approx 45{,}75\;\degC$.

\medskip

\textbf{Questions.}
\begin{itemize}[leftmargin=*]
  \item Quel pas donne le résultat le plus proche
        de la solution exacte ?
  \item Le pas $\Delta t = 5$ demande moins de calculs :
        que perd-on en contrepartie ?
\end{itemize}

\medskip

\textbf{Comparaison :}

\atrou[3.6cm]{%
\begin{itemize}
  \item $\Delta t = 2$ : $T \approx 42{,}94\;\degC$,
        écart $\approx 2{,}8\;\degC$.
  \item $\Delta t = 5$ : $T \approx 37{,}5\;\degC$,
        écart $\approx 8{,}3\;\degC$.
  \item Pas plus petit $\Rightarrow$ plus de calculs,
        mais meilleure précision.
\end{itemize}
}

% ====================================================================
\section*{7. Peut-on faire confiance au résultat ?}
% ====================================================================

On vient de calculer $T(10) \approx 42{,}94\;\degC$ avec
$\Delta t = 2$ et $T(10) \approx 37{,}5\;\degC$ avec $\Delta t = 5$,
alors que la valeur exacte est $45{,}75\;\degC$.
Ni l'un ni l'autre ne tombent juste.
D'où viennent ces écarts ?
Et si l'on confie le calcul à un ordinateur,
peut-on faire une confiance absolue au résultat affiché ?

\medskip

\textbf{Quiz.}
\begin{itemize}[leftmargin=*]
  \item Que vaut $0{,}1 + 0{,}2$ sur un ordinateur ?
  \item $\sqrt{2}^{\,2}$ vaut-il exactement $2$ ?
  \item Une petite erreur répétée peut-elle s'accumuler ?
\end{itemize}

\medskip

\textbf{Trois sources d'erreur :}

\atrou[6cm]{%
\begin{enumerate}
  \item \textbf{Erreur de modèle} :
    le café réel $\neq$ modèle de Newton.
  \item \textbf{Erreur de méthode} :
    Euler remplace la dérivée par
    $\dfrac{T_{n+1}-T_n}{\Delta t}$.
  \item \textbf{Erreur d'arrondi} :
    la machine stocke des nombres en précision finie.
\end{enumerate}
}

\textbf{Mesurer un écart :}

\atrou[4.4cm]{%
\begin{itemize}
  \item \textbf{Erreur absolue} :
    $\bigl|T_{\mathrm{exact}} - T_{\mathrm{calculé}}\bigr|$.
  \item \textbf{Erreur relative} :
    $\dfrac{\bigl|T_{\mathrm{exact}}
      - T_{\mathrm{calculé}}\bigr|}
    {\bigl|T_{\mathrm{exact}}\bigr|}$.
  \item \textbf{Chiffres significatifs} :
    nombre de décimales fiables du résultat.
\end{itemize}
}

% ====================================================================
\section*{8. Que se passe-t-il si le pas est trop grand ?}
% ====================================================================

On a vu qu'un pas plus petit donne un meilleur résultat,
mais coûte plus de calculs.
Peut-on aller dans l'autre sens et prendre un pas très grand ?
Poussons l'expérience à l'extrême.

\medskip

Avec $k = 0{,}1$, $T_0 = 90\;\degC$, $T_{\mathrm{amb}} = 20\;\degC$
et une pente initiale de $-7$ :

\begin{center}
\renewcommand{\arraystretch}{1.4}
\begin{tabular}{ccc}
  \hline
  $\Delta t$ & $T_1$ & Commentaire \\
  \hline
  $2$ & $76\;\degC$ & raisonnable \\
  $10$ & $20\;\degC$ & équilibre en une étape \\
  $15$ & $-15\;\degC$ & température négative ! \\
  \hline
\end{tabular}
\end{center}

\medskip

\atrou[4.4cm]{%
\begin{itemize}
  \item Un pas trop grand peut produire un résultat
        \textbf{physiquement absurde}.
  \item Plus le pas est petit, plus le calcul est fidèle
        --- mais il coûte plus cher.
  \item Le choix du pas est un \textbf{compromis
        entre précision et coût}.
\end{itemize}
}

\medskip

\textit{Au CM2, on formalisera ces notions de précision,
stabilité et convergence, et on verra comment la même
logique de discrétisation s'applique à des problèmes
plus complexes --- par exemple la diffusion de chaleur
dans une paroi.}

% ====================================================================
\section*{9. Bilan}
% ====================================================================

Trois idées à retenir :

\atrou[7cm]{%
\begin{enumerate}
  \item Un calcul numérique est une \textbf{approximation}.
  \item Une EDO peut être transformée en procédure de calcul
        par \textbf{discrétisation} (méthode d'Euler) :
        \[
          y_{n+1} = y_n + \Delta t\,f(t_n,\,y_n).
        \]
  \item Le choix du pas $\Delta t$ est un
        \textbf{compromis entre précision et coût}.
\end{enumerate}
}

\end{document}
